\documentclass[11pt]{article}
\usepackage[margin=1in]{geometry}
\usepackage[T1]{fontenc}
\usepackage[utf8]{inputenc}
\usepackage{lmodern}
\usepackage{enumitem}
\usepackage{xurl}
\urlstyle{tt}

\setlength{\parindent}{0pt}
\setlength{\parskip}{0.7\baselineskip}

\title{Prijedlog pravila za Python MIDI Enhancer}
\author{}
\date{}

\begin{document}
\maketitle

\section{Cilj projekta}

Python MIDI Enhancer učitava, analizira i poboljšava MIDI datoteke. Program mora
sačuvati glazbenu namjeru originala, a svaka izmjena mora biti razumljiva,
mjerljiva i reverzibilna.

\section{Korist, organizacija i način rada sustava}

\subsection{Glavne koristi}

Najveća korist sustava nije samo automatska promjena velocityja, nego
razumijevanje uloge svakog MIDI događaja prije donošenja odluke. Time se
postižu sljedeće koristi:

\begin{enumerate}[label=B\arabic*.]
    \item \textbf{Zaštita glazbene namjere.} Solo, riff, podloga, bass, Guitar
    Mode i pojedinačni drum element ne primaju isto pravilo samo zato što se
    nalaze u istom Styleu ili koriste isti Sound.

    \item \textbf{Točnost za Pa800.} Instrument se povezuje punom adresom
    \texttt{CC00.CC32.PC}, dok se RX, Guitar Mode, Drum Kit i drugi posebni
    mehanizmi obrađuju prema vlastitim pravilima.

    \item \textbf{Korištenje tvorničkog iskustva.} Factory Styles daju
    empirijske primjere kako je Korg koristio instrument u Intro, Variation,
    Fill, Break i Ending kontekstu. Ti podaci dopunjuju priručnik, ali ga ne
    zamjenjuju.

    \item \textbf{Odvojeni profili iste zvučne boje.} Isti instrument može
    imati različit radni velocity, trajanje i gustoću kao solo, akordska
    podloga, riff ili fixed-note fraza. Sustav te profile ne miješa.

    \item \textbf{Objašnjiv rezultat.} Svaka klasifikacija i izmjena navodi
    korišteno pravilo, ulazna mjerenja, razinu sigurnosti i izvor dokaza.

    \item \textbf{Sigurnost i reverzibilnost.} Original ostaje netaknut, a
    izlaz se ponovno parsira i uspoređuje događaj po događaj prije prihvaćanja.

    \item \textbf{Postupan razvoj.} Loader, klasifikator, pojedini Engine i
    verifier mogu se završiti i testirati zasebno. Pogreška u jednom modulu ne
    mora ugroziti cijelu aplikaciju.

    \item \textbf{Proširivost.} Novi instrument, Style izvor ili budući uređaj
    dodaje se profilom i adapterom, bez prepisivanja cijelog optimizera.
\end{enumerate}

\subsection{Planirana organizacija}

Sustav se organizira u slojeve s jasno ograničenom odgovornošću:

\begin{description}[style=nextline]
    \item[Core Loader i MIDI parser]
    Čitaju bajtove, provjeravaju strukturu, stvaraju kontrolni zbroj i grade
    nepromjenjivi model događaja. Ne donose glazbene odluke.

    \item[Source adapteri]
    Tumače strukturu konkretnog izvora, primjerice obični SMF, Pa800 Style SMF
    ili Style Works kolekciju. Adapter određuje Style Element, CV, valjani broj
    taktova i ulogu kanala, ali ne optimizira sadržaj.

    \item[Pa800 registar]
    Sadrži službenu adresu, naziv, vrstu i stranice priručnika za svaki Sound i
    Drum Kit. To je identitetski sloj, a ne skup preporučenih izmjena.

    \item[Context Classifier]
    Određuje je li zapis Drum, Perc, Bass, chordal accompaniment, line/riff,
    solo-kandidat, fixed Intro/Ending, Guitar/RX obrazac ili nepoznat slučaj.

    \item[Zajednički Instrument Engine]
    Računa univerzalna mjerenja: note, velocity, trajanje, raspon, polifoniju,
    gustoću, kontrolere, artikulacije i vremenski položaj.

    \item[Family Engine]
    Dodaje pravila obitelji kao što su bass, piano, strings, brass, reed,
    guitar, organ, synth i drugi tipovi instrumenata. Akustična, električna,
    elektromehanička i sintetizirana izvedba ostaju odvojene kada je razlika
    potvrđena.

    \item[Specijalizirani Engine]
    Obrađuje ponašanje koje se ne može svesti na obične note: Drum Engine,
    Guitar Mode, RX/DNC artikulacije, NRPN i druge službeno potvrđene mehanizme.

    \item[Profil pojedinog instrumenta]
    Za svaki \texttt{CC00.CC32.PC} čuva službeni identitet i odvojene
    empirijske profile za solo, chordal, riff, fixed i druge potvrđene načine
    korištenja. Profil je podatak, a ne zasebna kopija programskog koda.

    \item[Policy Engine]
    Uspoređuje analizu s dopuštenim pravilima i odlučuje smije li proizvesti
    prijedlog. Primjerice, 35\% dopušta samo potvrđenoj podlozi, a 50\%
    potvrđenoj sekundarnoj perkusiji.

    \item[Change Engine]
    Na radnu kopiju primjenjuje samo korisnički potvrđene prijedloge. Ne smije
    samostalno mijenjati instrument, tonalitet, note ili strukturu.

    \item[Verifier]
    Ponovno parsira rezultat, provjerava Note On/Off parove, strukturu,
    kontrolere i sve neizmijenjene događaje te blokira neispravan izlaz.

    \item[Report i GUI]
    Prikazuju činjenice, procjene, dokaze, upozorenja, prijedloge i razliku
    između originala i radne kopije. GUI ne sadrži skrivena glazbena pravila.
\end{description}

\subsection{Kako će obrada raditi}

Jedno pokretanje prolazi kroz strogo definiran tok:

\begin{enumerate}[label=W\arabic*.]
    \item Korisnik odabire MIDI, mapu ili ZIP. Loader bilježi metapodatke i
    SHA-256 te original otvara samo za čitanje.

    \item Parser provjerava MIDI strukturu i pretvara događaje u nepromjenjivi
    zajednički model, uz očuvanje sirovih bajtova i izvornog redoslijeda.

    \item Source adapter prepoznaje vrstu izvora. Kod Style Worksa razdvaja
    Style, Element, CV i uloge te ograničava analizu na deklarirani broj
    taktova.

    \item Instrument Resolver prati Bank Select i Program Change po kanalu i
    vremenu te povezuje puni \texttt{CC00.CC32.PC} sa službenim registrom.

    \item Context Classifier zasebno klasificira svaki Element/CV/track. Bez
    potvrđene ili dovoljno pouzdane uloge obrada ostaje read-only.

    \item Odgovarajući zajednički, obiteljski i specijalizirani Engine računaju
    profil. Guitar Mode ide Guitar Engineu, Drum Kit Drum Engineu, a normalni
    bass Bass Family Engineu.

    \item Profil se uspoređuje samo s odgovarajućim Factory Style kontekstom.
    Solo se uspoređuje sa solo primjerima, a rhythm guitar s rhythm-guitar
    primjerima; statistike različitih uloga ne miješaju se.

    \item Analyze prikazuje izmjerene činjenice. Suggest dodaje moguće izmjene
    s formulom, rizikom i pravilom. Nijedan prijedlog još ne mijenja MIDI.

    \item Korisnik pojedinačno ili jasno ograničenom grupom potvrđuje željene
    prijedloge. Change Engine ih primjenjuje na zasebnu radnu kopiju.

    \item Verifier uspoređuje radnu kopiju s originalnim modelom. Ako postoji
    neočekivana razlika, spremanje se blokira i rezultat dobiva status
    \texttt{FAIL}.

    \item Tek nakon uspješne provjere stvaraju se nova MIDI datoteka, JSON
    izvještaj, dnevnik promjena i SHA-256 izlaza.
\end{enumerate}

Ovakva organizacija znači da nijedan Engine ne mora znati kako se čita ZIP,
crta GUI ili sprema datoteka. Svaki modul prima provjerene podatke, izvršava
jedan ograničen zadatak i vraća rezultat s dokazom.

\section{Osnovna pravila}

\begin{enumerate}[label=R\arabic*.]
    \item Originalna MIDI datoteka nikada se ne mijenja izravno.
    Rezultat se sprema kao nova datoteka.

    \item Analiza se uvijek izvršava prije bilo kakvog poboljšavanja.

    \item Program mora jasno razdvojiti činjenice pronađene u MIDI datoteci od
    procjena i prijedloga koje generira enhancer.

    \item Korisnik mora moći uključiti ili isključiti svaku vrstu poboljšanja.

    \item Zadane postavke moraju biti konzervativne i ne smiju značajno
    promijeniti melodiju, harmoniju, ritam ili dinamiku.

    \item Svaka promjena mora biti zabilježena u izvještaju: vrsta promjene,
    lokacija, stara vrijednost i nova vrijednost.

    \item Prije spremanja program mora provjeriti rezultat. Mora prijaviti
    neispravne, izgubljene i nezatvorene note.
\end{enumerate}

\section{Jedini izvor istine za Korg Pa800}

Normativni izvori za Pa800-specifične funkcije službeni su Korgovi priručnici
\emph{KORG Pa800 2.0 User's Manual (E12)} i \emph{KORG Pa800 2.0 Advanced Edit
Manual (E5)}. Lokalno su spremljeni kao
\path{prism-uploads/Pa800-201UM-ENG.pdf} i
\path{prism-uploads/Pa800_AE_E.pdf}. Forumi,
videozapisi, drugi modeli Korg uređaja i programske pretpostavke ne mogu biti
dokaz ponašanja Pa800 uređaja.

User's Manual ima 344 PDF stranice i SHA-256 kontrolni zbroj
\path{b7df16fb0113d998a6411ea1c28971707021da8bc0b728aff98eafcceb57c71b}.
Advanced Edit Manual ima 124 PDF stranice i SHA-256 kontrolni zbroj
\path{6bcda56658a89eda31e9fe62a2b744df5cecaaf478e050d3256ad6440c0d413e}.
Za popis tvorničkih Soundova i Drum Kitova normativan je User's Manual, dok se
Advanced Edit Manual koristi za detalje uređivanja i ponašanja zvuka.

Svako tehničko pravilo vezano uz Pa800 mora u dnevniku navesti točne stranice
priručnika na kojima je potvrđeno. Ako priručnik ne daje odgovor, funkcija se
označava kao nepotvrđena i ne implementira se kao činjenica.

\section{Modularni način razvoja}

Projekt se od sada razvija strogo modul po modul. Za svaki modul vrijedi
sljedeći redoslijed:

\begin{enumerate}[label=M\arabic*.]
    \item definirati jednu ograničenu funkciju i kriterij uspjeha;
    \item pronaći i zapisati relevantne stranice Pa800 priručnika;
    \item implementirati samo definirani opseg;
    \item izvršiti automatski test i sačuvati naziv testne naredbe;
    \item zapisati koje su datoteke i funkcije promijenjene;
    \item dodijeliti status \texttt{PASS}, \texttt{FAIL} ili
    \texttt{BLOCKED};
    \item prijeći na ovisni modul samo kada su svi njegovi deklarirani
    preduvjeti \texttt{PASS}; nepovezani \texttt{BLOCKED} ili \texttt{FAIL}
    modul ne zaustavlja neovisnu razvojnu granu.
\end{enumerate}

Trajna evidencija nalazi se u datoteci \texttt{MODULE\_LOG.md}. Status
\texttt{PASS} znači da je funkcija implementirana, test izvršen i dokaz iz
priručnika naveden kada je funkcija Pa800-specifična. Status \texttt{FAIL}
znači da test nije prošao, dok \texttt{BLOCKED} označava nedostajući službeni
izvor, biblioteku, grafičko okruženje ili drugi obavezni preduvjet.

\section{Zatvoreni sustav pravila i dokazivanja}

Nijedna odluka aplikacije ne smije nastati iz prešutne pretpostavke. Svaki
rezultat, upozorenje, prijedlog ili izmjena mora navesti pravilo koje ga je
proizvelo, ulazne podatke na kojima se pravilo temelji i razinu sigurnosti.
Obavezni statusi svake tvrdnje su:

\begin{description}[style=nextline]
    \item[\texttt{CONFIRMED}] Tvrdnja je izravno zapisana u MIDI događaju,
    službenom Pa800 priručniku ili drugom unaprijed odobrenom normativnom
    izvoru.
    \item[\texttt{DERIVED}] Vrijednost je deterministički izračunata iz
    potvrđenih događaja, primjerice broj nota, trajanje ili polifonija.
    \item[\texttt{INFERRED}] Rezultat je heuristička procjena. Mora sadržavati
    pravilo, korištena mjerenja i postotak pouzdanosti.
    \item[\texttt{CONFLICT}] Dva ili više izvora daju različite rezultate.
    Nijedan se ne proglašava pobjednikom bez unaprijed definiranog pravila
    prioriteta.
    \item[\texttt{UNKNOWN}] Nema dovoljno podataka za zaključak.
    \item[\texttt{UNSUPPORTED}] Događaj je prepoznat kao valjan, ali ga
    trenutni modul ne zna protumačiti.
\end{description}

Za statuse \texttt{CONFLICT}, \texttt{UNKNOWN} i \texttt{UNSUPPORTED} vrijedi
obavezno završno pravilo: sačuvati izvorne bajtove, prijaviti razlog i ne
primijeniti izmjenu. Nedostatak pravila nikada ne znači dopuštenje za zadanu
vrijednost, General MIDI zamjenu, brisanje događaja ili nagađanje instrumenta.

\section{Pravila učitavanja i analize}

Program u početnoj fazi mora prepoznati:

\begin{itemize}
    \item MIDI format, broj trackova i rezoluciju u tickovima;
    \item trajanje, tempo i promjene tempa;
    \item takt i promjene oznake takta;
    \item note, njihov početak, trajanje i velocity;
    \item kanale, programe, instrumente i percussion kanal;
    \item raspon tonova i polifoniju svakog tracka;
    \item kontrolne poruke, pitch bend i sustain pedal;
    \item moguće probleme u strukturi datoteke.
\end{itemize}

Nepoznati događaji ne smiju se automatski brisati. Ako ih program ne može
obraditi, mora ih sačuvati ili jasno upozoriti korisnika.

\subsection{Stroga pravila učitavanja}

Učitavanje je zaseban modul koji ne optimizira i ne ispravlja sadržaj. Za svaki
ulaz vrijede sljedeća pravila:

\begin{enumerate}[label=L\arabic*.]
    \item Prije parsiranja bilježe se izvorni naziv, relativna putanja,
    veličina, vrijeme izmjene i SHA-256 kontrolni zbroj.

    \item Ekstenzija \texttt{.mid} ili \texttt{.midi} nije dokaz valjanosti.
    Datoteka se prihvaća kao Standard MIDI File tek nakon provjere zaglavlja
    \texttt{MThd}, duljine zaglavlja i svih \texttt{MTrk} blokova.

    \item Provjeravaju se MIDI format 0, 1 ili 2, deklarirani broj trackova,
    division, granice svakog bloka, variable-length vrijednosti, running
    status i End-of-Track događaj.

    \item SMPTE division ne smije se tumačiti kao PPQ. Ako modul nema podršku
    za SMPTE vrijeme, datoteka se učitava samo za očuvanje i dobiva status
    \texttt{UNSUPPORTED} za glazbenu vremensku analizu.

    \item Svaki događaj dobiva fizički track, MIDI kanal ako postoji,
    apsolutni tick, izvorni redni broj na tom ticku, delta vrijeme i sirove
    bajtove. Time se kasnije može dokazati da struktura nije promijenjena.

    \item Running status mora se razriješiti za analizu, ali se pri spremanju
    ne smije proizvoljno prepisati ako je cilj bajtovno očuvanje neizmijenjenih
    događaja.

    \item Note On s velocityjem nula analizira se kao Note Off, ali se njegov
    izvorni oblik pamti i ne pretvara automatski u drugi statusni bajt.

    \item Nepotpuna nota, događaj izvan granice tracka, nevažeći VLQ ili
    nedostajući End-of-Track prijavljuje se s točnom lokacijom. Parser ne smije
    nastaviti preko nepoznate granice kao da je sadržaj ispravan.

    \item SysEx, meta-događaji, RPN, NRPN i nepoznati kontroleri čuvaju se u
    izvornom redoslijedu čak i kada ih analiza ne koristi.

    \item ZIP ili mapa obrađuje se kao skup zasebnih izvora. Ne smije se
    mijenjati arhiv, redoslijed datoteka ni nazivi bez izričitog izvoza u novu
    lokaciju.

    \item Dvije datoteke jednakog naziva ne smatraju se jednakima bez jednakog
    kontrolnog zbroja. Duplikati se bilježe, ali se ne brišu automatski.

    \item Nakon učitavanja ponovno se računa SHA-256 originala. Ako se razlikuje
    od početnog zbroja, obrada se prekida statusom \texttt{FAIL}.
\end{enumerate}

\subsection{Pravila identiteta instrumenta}

Instrument se povezuje prvenstveno prema stvarnom MIDI stanju kanala, a ne
prema nazivu tracka. Identitet se izračunava sljedećim redoslijedom:

\begin{enumerate}[label=I\arabic*.]
    \item \texttt{CC0} postavlja Bank Select MSB, \texttt{CC32} postavlja Bank
    Select LSB, a Program Change aktivira program u trenutno odabranoj banci.
    Puni identifikator prikazuje se kao \texttt{CC00.CC32.PC}.

    \item Bank Select stanje prati se zasebno za svaki MIDI kanal i kroz
    vremenski slijed događaja. Kasniji Program Change ne smije koristiti
    bankovne vrijednosti drugog kanala.

    \item Identitet vrijedi od ticka Program Change događaja do sljedeće
    potvrđene promjene. Jedan track zato može imati više vremenskih segmenata i
    više instrumenata.

    \item Kada su sva tri podatka prisutna, adresa se traži u službenom Pa800
    registru. Potpuno podudaranje dobiva status \texttt{CONFIRMED}.

    \item Tekstualni naziv Sounda ili instrumenta služi kao neovisna kontrola.
    Ako se naziv i službena adresa ne slažu, rezultat je \texttt{CONFLICT};
    adresa se ne prepisuje prema tekstu niti se tekst briše.

    \item Ako nedostaje \texttt{CC0} ili \texttt{CC32}, vrijednost se može
    naslijediti samo iz ranijeg događaja istoga kanala u istom MIDI toku.
    Vrijednost se ne prenosi iz drugog CV-a, Style Elementa, datoteke ili
    susjednog tracka.

    \item Banka \texttt{0.0} ne pretpostavlja se samo zato što kontroleri
    nedostaju. GM zadana banka dopuštena je jedino kada je GM/GM2 način
    potvrđen odgovarajućim SysExom, dokumentiranim formatom izvora ili izričitom
    korisničkom postavkom.

    \item Sam Program Change, naziv tracka ili glazbeni raspon nisu dovoljni za
    potvrdu Pa800 Sounda. Takav rezultat ostaje \texttt{INFERRED} ili
    \texttt{UNKNOWN}.

    \item Svaki Style Element, CV i uloga Drum, Perc, Bass ili Acc1--Acc5 ima
    vlastiti zapis instrumenta. Instrument se ne prenosi između njih samo zato
    što imaju jednak naziv uloge.

    \item User Sound i User Drum Kit adresa potvrđuje lokaciju, ali ne i sadržaj
    korisničkog slota. Naziv i stvarni multisample ostaju \texttt{UNKNOWN} bez
    podataka iz konkretnog uređaja ili SET-a.

    \item Remap adresa čuva i traženu adresu i službenu ciljnu adresu. Izvještaj
    ne smije sakriti da je uređaj izvršio remap.

    \item Drum Kit je identitet kompleta, dok je svaka korištena MIDI nota
    zaseban drum element. Element dobiva naziv samo iz mappinga potvrđenog za
    taj konkretni kit.
\end{enumerate}

\subsection{Odvajanje obitelji instrumenata i Engine profila}

Klasifikator najprije odvaja \emph{pitched instrument}, \emph{Drum Kit},
\emph{percussion kit}, \emph{sound effect} i \emph{unknown}. Bubanj se ne
obrađuje kroz Engine melodijskih instrumenata. Za melodijske instrumente Engine
odvojeno bilježi akustičnu, električnu, elektromehaničku i sintetiziranu
obitelj. Ako službeni naziv ili adresa ne potvrđuju obitelj, klasifikacija ostaje
\texttt{INFERRED}.

Svaki profil melodijskog instrumenta mora sadržavati puni službeni identitet,
obitelj, fizički ili očekivani tonski raspon, potvrđeni raspon korišten u
uzorku, minimalni i maksimalni velocity, radni velocity interval, distribuciju
i krivulju velocityja, tipičnu polifoniju, artikulacije, trajanje, gustoću,
pedale, kontrolere, pitch bend, ulogu solo ili podloga te izvore svakog pravila.
Vrijednosti dobivene iz Factory Styles označavaju se kao empirijski profil, a
ne kao fizička granica stvarnog instrumenta.

Empirijski podaci ne smiju se spojiti u jedan zajednički profil instrumenta.
Za svaku službenu adresu \texttt{CC00.CC32.PC} vode se odvojeni kontekstualni
profili: \texttt{ACCOMP\_CHORDAL}, \texttt{ACCOMP\_LINE\_RIFF},
\texttt{SOLO\_CANDIDATE}, \texttt{FIXED\_INTRO\_ENDING},
\texttt{GUITAR\_RX\_PATTERN} i \texttt{UNKNOWN}. Velocity, raspon, trajanje i
gustoća solo zapisa nikada se ne koriste za izračun podloge, niti se statistika
podloge primjenjuje na solo zapis.

Drum Engine za svaku notu zasebno čuva kit adresu, broj note, potvrđeni naziv
elementa, glavnu ili sekundarnu ulogu, velocity statistiku, položaj u taktu,
učestalost, flam ili simultani udar, choke grupu ako je dokazana, NRPN stanje i
mogući RX/DNC ili velocity-switch rizik. Nepoznat mapping ne smije se zamijeniti
General MIDI nazivom.

\subsection{Poseban analizator Pa800 Factory Styles}

Factory Style korpus analizira se odvojeno od optimizera. Njegov zadatak je
izvući dokazive obrasce, a ne mijenjati izvezene MIDI datoteke.

Ovaj analizator nema funkciju postavljanja zvuka. Ne smije dodavati niti
mijenjati Bank Select, Program Change, Sound, Style Performance ili Original
Style Sounds postavku. Adresa \texttt{CC00.CC32.PC} koristi se isključivo kao
ključ kojim se opaženi obrazac povezuje s već registriranim instrumentom.
Factory Styles dopunjuju ono što User's Manual ne propisuje kao fiksnu
vrijednost, primjerice stvarne velocity distribucije, gustoću, fraziranje i
način korištenja instrumenta u različitim Style kontekstima.

\begin{enumerate}[label=S\arabic*.]
    \item Mapa predstavlja kolekciju, podmapa kandidata za Style, a svaka MIDI
    datoteka kandidata za jedan Style Element. Naziv datoteke je trag, ali se
    potvrđuje sadržajem i tekstualnim meta-događajima.

    \item Obavezno se razlikuju Variation 1--4, Intro 1--3, Fill 1--2,
    Break/Fill 3 i Ending 1--3. Nijedan element ne smije se spojiti s drugim u
    zajedničku velocity statistiku bez oznake izvora.

    \item Svaki element se dijeli na CV1--CV6 kada postoje. Variation može imati
    do šest CV-ova, a ostali elementi do dva. Prazni CV placeholder nije isto
    što i nedostajući ili oštećeni CV.

    \item Svaki CV mora zasebno mapirati Drum, Perc, Bass i Acc1--Acc5. Službeni
    kanalni raspored 10, 11, 9 i 12--16 koristi se kao dokaz uloge samo za
    potvrđeni Pa800 Style izvor.

    \item Style Works Format 1 izvoz obrađuje poseban read-only adapter. Ne
    proglašava se izvornim Pa800 Format 0 Style MIDI-jem i ne dobiva izmišljene
    Korg markere.

    \item Deklarirani broj taktova uspoređuje se s vremenskim potpisom i zadnjim
    sadržajnim događajem. Umjetno udaljen End-of-Track ne smije povećati
    glazbeno trajanje elementa.

    \item Style Works track može iza deklarirane granice sadržavati zaostale
    događaje i tekst drugih Style Elemenata. Valjani prozor završava na
    \texttt{broj taktova × trajanje takta}; događaji iza njega ne ulaze u
    statistiku aktivnog elementa, ali se čuvaju i prijavljuju.

    \item Za svaki element i CV mjere se: note, velocity distribucija i
    percentili, pozicija naglasaka, gustoća, sinkopa, swing, trajanja,
    polifonija, instrumenti, kontroleri, raspon, ponavljanje uzorka i tišina.

    \item Analiza prijelaza uspoređuje Variation--Fill--Variation,
    Intro--Variation i Variation--Ending. Fill se ne ocjenjuje kao samostalna
    petlja bez konteksta ulazne i izlazne varijacije.

    \item Analiza Acc trackova mora razlikovati fiksne note, transponirajuće
    akorde, riff, arpeggio, pad i solo frazu. Bez službenih NTT/NTR, Key/Chord i
    Chord Table podataka rezultat ostaje heuristički.

    \item Parametri koji nisu prisutni u Style Works MIDI izvozu ne smiju se
    rekonstruirati kao činjenice. Bilježe se u popisu nedostajućih podataka za
    eventualni SET, uređaj ili ručni unos.

    \item Nepotpuni Style ostaje analiziran kao nepotpun skup. Program ne smije
    kopirati element iz drugog Stylea da bi umjetno zadovoljio očekivanih 13
    elemenata.
\end{enumerate}

\subsection{Stroga razlika između ritma, podloge i solo tracka}

Pojam \emph{Style track} opisuje mjesto unutar Style strukture, ali ne dokazuje
da je sadržaj ritmička podloga. User's Manual na tiskanoj stranici 134
(PDF stranica 138) izričito dopušta da Acc track sadrži melodijski ili
harmonijski obrazac. Na tiskanoj stranici 116 (PDF stranica 120) navodi se da su
Parallel NTT tablice prikladne melodijskim dijelovima, Fixed tablice akordskim
trackovima, a Intro 1 i Ending 1 u tvorničkim Styleovima standardno koriste
No Transpose i često sadrže snimljenu progresiju. Style Works izvoz ne sadrži
NTT, NTR, Key/Chord ni Chord Table; njihovo odsustvo potvrđuje i postupak SMF
uvoza na tiskanoj stranici 135 (PDF stranica 139).

Zbog toga se klasifikacija izvršava sljedećim redoslijedom:

\begin{enumerate}[label=T\arabic*.]
    \item \texttt{RHYTHM\_DRUM} potvrđuje se samo za Drum track i
    \texttt{RHYTHM\_PERC} samo za Perc track potvrđen Pa800 kanalnim rasporedom
    ili drugim izričitim metapodatkom. Oni se obrađuju kroz Drum Engine, ne kroz
    melodijski Instrument Engine.

    \item Bass track dobiva status \texttt{BASS\_ACCOMP}. Prema priručniku
    njegova je Style funkcija pratiti korijen prepoznatog akorda. Pokretna bass
    linija ne proglašava se solo dionicom samo zato što je monofona.

    \item Prije melodijske klasifikacije moraju se izdvojiti Guitar Mode
    kontrolne note, strumming naredbe, power-chord tipke i RX noise/articulation
    note. Ako Track Type i mapping nisu dostupni, zapis dobiva
    \texttt{GUITAR\_RX\_PATTERN} ili \texttt{UNKNOWN}; velik tonski raspon tada
    nije dokaz solo sviranja.

    \item Acc zapis je \texttt{ACCOMP\_CHORDAL} kada najmanje 30\% početaka
    sadrži dvije ili više istodobnih nota ili kada jedan početak sadrži najmanje
    tri note. Ovo opisuje teksturu; ne dokazuje skriveni NTT Track Type.

    \item Acc zapis je konzervativni \texttt{SOLO\_CANDIDATE} samo kada nakon
    uklanjanja potvrđenih kontrolnih/artikulacijskih nota ima najmanje šest
    nota, najmanje 95\% monofonih početaka, najviše dvije istodobne i dvije
    preklopljene note, najmanje pet različitih visina, raspon od najmanje devet
    polutonova i prosječno melodijsko kretanje od najmanje 1,5 polutona.

    \item Monofoni Acc zapis koji pokazuje ponavljajući motiv, arpeggio,
    ostinato, kratko ritmičko sjeckanje ili ne zadovoljava sve solo uvjete dobiva
    \texttt{ACCOMP\_LINE\_RIFF}. Ne smije se automatski proglasiti ni podlogom
    ni solom.

    \item Melodijski ili akordski sadržaj Intro 1 i Ending 1 dodatno dobiva
    oznaku \path{FIXED_INTRO_ENDING_CANDIDATE}. Oznaka postaje
    \texttt{CONFIRMED} samo ako SET, uređaj ili drugi službeni izvor potvrdi
    No Transpose/NTT postavku.

    \item Fill, Break, Intro, Ending i Variation analiziraju se odvojeno.
    Kratka prijelazna fraza ne ulazi u radni velocity profil kontinuirane
    pratnje samo zato što koristi isti instrument i isti Acc broj.

    \item Strings, brass, guitar, piano ili bilo koja druga obitelj ne dobiva
    ulogu prema nazivu instrumenta. Strings može biti akordska podloga,
    monofona kontramelodija ili solo; odluku donosi sadržaj pojedinog
    Elementa/CV-a.

    \item Samo \texttt{ACCOMP\_CHORDAL} koji je dodatno potvrđen kao Lower,
    pad ili chord podloga može koristiti pravilo smanjenja od 35\%.
    Statusi \path{SOLO_CANDIDATE}, \path{ACCOMP_LINE_RIFF},
    \path{FIXED_INTRO_ENDING_CANDIDATE}, \path{GUITAR_RX_PATTERN} i
    \path{UNKNOWN} ne mijenjaju velocity automatski.

    \item Korisnik može ručno potvrditi ili ispraviti ulogu, ali izvještaj mora
    sačuvati izvorni automatski rezultat, ručnu odluku i vrijeme odluke.
\end{enumerate}

\subsection{Rezultat analize učitanog Factory Style korpusa}

Arhiv \path{prism-uploads/Split Factory Styles.zip} koristi se kao empirijski,
read-only izvor i ima SHA-256
\path{ffb95cc191ba2bfc0c9b2b09494adaeb61c61bcda0fa4a9e4d83302e281d8f5e}.
Sadrži 3.211 valjanih MIDI datoteka i 252 prepoznata Style skupa. Nakon
ograničavanja svakog elementa na deklarirani broj taktova pronađeno je 29.639
aktivnih track zapisa.

Među 14.162 aktivna Acc zapisa s najmanje šest nota, početna verzija gore
navedenih pravila daje 6.554 \texttt{ACCOMP\_CHORDAL} zapisa (46,28\%), 3.362
\texttt{ACCOMP\_LINE\_RIFF} zapisa (23,74\%), 1.143 konzervativna
\texttt{SOLO\_CANDIDATE} zapisa (8,07\%), 2.449
\texttt{GUITAR\_RX\_PATTERN} zapisa (17,29\%) i 654 neriješena zapisa
(4,62\%). Ovo su rezultati determinističke klasifikacije, a ne tvorničke NTT
postavke.

Solo-kandidati nisu ograničeni na Intro. Pronađena su 303 među 1.840 analiziranih
Acc zapisa u Intro 1, 262 među 1.929 u Ending 1 te ukupno 379 među 5.956 Acc
zapisa u Variation 1--4. Time je dokazano da aplikacija ne smije primijeniti
jedan zajednički \emph{Style velocity} profil na sve Acc trackove.

\subsection{Minimalni rezultat analize}

Analiza mora proizvesti hijerarhijski rezultat koji redom obuhvaća kolekciju,
Style, Element, CV, ulogu, kanal, instrument i notu. Za svaku razinu navodi
izvor, status dokaza, upozorenja i mjerenja. Samo prikaz ukupnog minimuma,
maksimuma i prosjeka velocityja nije dovoljan kriterij završene analize.

\section{Pravila poboljšavanja}

\subsection{Timing i kvantizacija}

Kvantizacija mora imati podesivu jačinu. Potpuna kvantizacija nije zadana jer
može ukloniti prirodan osjećaj izvedbe. Note se ne smiju pomaknuti preko granice
koja bi promijenila njihov redoslijed ili proizvela neželjeno preklapanje.

\subsection{Dinamika}

Velocity se mijenja unutar dopuštenog raspona od 1 do 127. Enhancer mora
sačuvati relativne naglaske i ne smije sve note izjednačiti na istu vrijednost.

\subsection{Relativna dinamika podloga i sekundarne perkusije}

Smanjenje dinamike uvijek je relativno prema izvornoj velocity vrijednosti, a
ne postavljanje svih nota na fiksni broj. Time se čuvaju naglasci izvođača i
odnos tihih i snažnih udara.

\begin{enumerate}[label=V\arabic*.]
    \item Track potvrđen kao \emph{Lower}, chordal accompaniment ili harmonijska
    podloga smanjuje se za 35\%, odnosno nova vrijednost računa se kao
    \texttt{round(velocity * 0.65)}.

    \item Strings track potvrđen kao solo ili lead ne smanjuje se. Strings se
    smanjuje 35\% samo kada je potvrđen kao podloga, pad ili akordska pratnja.

    \item Sekundarna perkusija na MIDI kanalu 10 smanjuje se za 50\%, odnosno
    nova vrijednost računa se kao \texttt{round(velocity * 0.50)}.

    \item Kick, bass drum, snare, side stick, clap, hi-hat, tomovi i glavni
    cymbal elementi smatraju se glavnim drum elementima i ovo pravilo ih ne
    smanjuje.

    \item Conga, bongo, timbale, agogo, cabasa, maracas, whistle, guiro,
    claves, wood block, cuica, triangle i srodni ukrasni elementi smatraju se
    sekundarnom perkusijom kada je njihov identitet potvrđen.

    \item Izračunata Note On velocity vrijednost nikada ne smije pasti ispod 1
    niti prijeći 127. Note On s velocityjem 0 ostaje Note Off semantika i ne
    pretvara se u zvučnu notu.

    \item General MIDI note mapping koristi se samo kada je potvrđen GM ili GM2
    Drum Kit. Kod RX, User, oriental, ethnic i drugih Pa800-specifičnih kitova
    nota se ne proglašava određenim elementom samo prema GM tablici.

    \item Ako smanjenje može prijeći tvornički velocity-switch prag i aktivirati
    drugi sample ili RX sloj, promjena se u načinu \emph{Suggest} samo prijavljuje.
    Primjenjuje se tek kada su pragovi poznati ili korisnik izričito potvrdi
    promjenu.

    \item Ako uloga tracka nije sigurno potvrđena, velocity se ne mijenja.
    Program prikazuje prijedlog i razloge klasifikacije.
\end{enumerate}

Izraz \emph{kanal 10} označava MIDI kanal 10, odnosno internu nultu vrijednost
kanala 9. Ne označava deseti fizički track u MIDI datoteci.

\subsection{Stroga NRPN pravila za Pa800 Drum parametre}

NRPN poruke tretiraju se kao povezani, redoslijedom osjetljivi skup događaja.
Za Pa800 Drum parametre službeni slijed čine \texttt{CC99} kao NRPN MSB,
\texttt{CC98} kao NRPN LSB i \texttt{CC6} kao Data Entry. Vrijednost
\texttt{CC98} određuje drum element, odnosno MIDI notu 0--127. Ako izvor
sadrži \texttt{CC38} ili dodatne Data Entry poruke, one se čuvaju bez
nagađanja njihova značenja.

Službeni Pa800 NRPN MSB brojevi za Drum parametre su:

\begin{description}[style=nextline]
    \item[20] filter cutoff;
    \item[21] filter resonance;
    \item[22] EG attack time;
    \item[23] EG decay time;
    \item[24] coarse tune;
    \item[25] fine tune;
    \item[26] volume;
    \item[28] panpot;
    \item[29] reverb, odnosno FX1 send;
    \item[30] modulation, odnosno FX2 send.
\end{description}

NRPN skup mora ostati na istom MIDI kanalu, u istom tracku i u izvornom
redoslijedu. Program ga ne smije razdvojiti, sortirati samo prema broju
kontrolera, premjestiti iza druge note ili ukloniti kao višak. Nepoznat ili
nepotpun NRPN slijed čuva se nepromijenjen i prijavljuje kao upozorenje.

\subsection{Očuvanje strukture i tekstualnih događaja}

Promjena velocityja ne smije mijenjati MIDI format, division, broj i redoslijed
trackova, apsolutni tick događaja, redoslijed događaja na jednakom ticku,
Program Change, Bank Select, SysEx, NRPN, RPN ni druge kontrolne poruke.

Svi tekstualni meta-događaji čuvaju se bajt-po-bajt kad god njihov sadržaj nije
izričito uređivan. To uključuje opći tekst, copyright, naziv tracka, naziv
instrumenta, lyrics, marker i cue point. Nepoznat meta-događaj ili nestandardno
kodiranje ne smije se automatski normalizirati, prevesti ili izbrisati.

Prije spremanja uspoređuju se originalna i nova datoteka. Izvan izričito
odobrenih velocity vrijednosti ne smije postojati nijedna sadržajna razlika.
Rezultat se uvijek sprema kao nova MIDI datoteka.

\subsection{Trajanje nota}

Program može ukloniti vrlo kratke slučajne note i ispraviti problematična
preklapanja samo iznad korisnički definiranog praga. Takve promjene moraju biti
označene kao potencijalno destruktivne.

\subsection{Melodija i harmonija}

Automatska promjena visine tona, akorda ili tonaliteta nije dopuštena u
osnovnom načinu rada. Takve funkcije pripadaju posebnom kreativnom načinu i
zahtijevaju izričitu potvrdu korisnika.

\section{Načini rada}

\begin{description}[style=nextline]
    \item[Analyze] Samo učitavanje i analiza, bez promjene MIDI sadržaja.
    \item[Suggest] Analiza i popis predloženih promjena, bez automatske primjene.
    \item[Enhance] Primjena samo odabranih i potvrđenih poboljšanja.
\end{description}

\section{Izlaz programa}

Svako pokretanje treba moći proizvesti:

\begin{itemize}
    \item čitljiv sažetak analize;
    \item strukturirani JSON izvještaj;
    \item novu MIDI datoteku, ako su promjene primijenjene;
    \item dnevnik svih upozorenja i izvršenih promjena.
\end{itemize}

Naziv izlazne datoteke treba sadržavati sufiks, primjerice
\texttt{naziv\_enhanced.mid}, kako bi original ostao netaknut.

\subsection{Pravila izvještaja i spremanja}

JSON izvještaj mora sadržavati verziju sheme, podatke o izvoru, kontrolne
zbrojeve, odabrani parser, konfiguraciju, hijerarhiju Style podataka, sve
potvrđene i procijenjene instrumente, upozorenja, prijedloge i dnevnik promjena.
Svaka izmjena mora imati stabilni ID pravila kako bi se rezultat mogao ponovno
izračunati i testirati.

Spremanje izmijenjenog MIDI-ja dopušteno je samo ako su ispunjene sve kapije:

\begin{enumerate}[label=E\arabic*.]
    \item originalni SHA-256 i dalje odgovara datoteci koja je učitana;
    \item nema neriješenog \texttt{BLOCKING} upozorenja;
    \item svaka izmjena pripada izričito uključenom i potvrđenom pravilu;
    \item rezultat se može ponovno parsirati bez nove strukturne greške;
    \item broj nota, Note On/Off parovi i vrijeme događaja odgovaraju očekivanoj
    vrsti izmjene;
    \item svi neizmijenjeni događaji jednaki su izvornim događajima;
    \item izlazni naziv nije jednak ulaznom nazivu u istoj lokaciji;
    \item izračunat je SHA-256 izlaza i spremljen u izvještaj.
\end{enumerate}

Ako ijedna kapija ne prođe, privremeni rezultat se ne proglašava dovršenim
MIDI-jem. Aplikacija čuva dijagnostički izvještaj i dodjeljuje status
\texttt{FAIL}.

\section{Pravila konfiguracije, testiranja i verzioniranja}

Svi pragovi, uključujući 35\% za Lower/chord podloge i 50\% za sekundarne
perkusije, moraju biti imenovane konfiguracijske vrijednosti s jedinicom,
dopuštenim rasponom, zadanom vrijednošću i izvorom pravila. Nevažeća vrijednost
ne smije se tiho ograničiti ili zamijeniti zadanom; konfiguracija se odbija uz
jasnu poruku.

Svaki modul mora imati najmanje:

\begin{itemize}
    \item unit test za svako pravilo i graničnu vrijednost;
    \item test valjanog, nepotpunog, konfliktnog i nepodržanog ulaza;
    \item round-trip test koji dokazuje očuvanje neizmijenjenih događaja;
    \item regression test za svaki pronađeni stvarni problem;
    \item golden izvještaj za najmanje jedan mali, ručno provjeren MIDI primjer;
    \item test determinističnosti ponovljenim pokretanjem iste analize;
    \item test da Analyze i Suggest ne mijenjaju izvorni kontrolni zbroj.
\end{itemize}

Promjena registra, JSON sheme, klasifikacijskog pravila ili formule zahtijeva
novu verziju odgovarajuće komponente. Stari izvještaj mora ostati čitljiv ili
dobiti jasan status nepodržane verzije; ne smije se tiho reinterpretirati prema
novim pravilima.

\section{Predloženi redoslijed razvoja}

\begin{enumerate}
    \item pouzdano učitavanje i analiza MIDI datoteke;
    \item validacija i izrada izvještaja;
    \item prikaz piano-rolla i statistike;
    \item prijedlozi za timing, velocity i trajanje nota;
    \item kontrolirana primjena poboljšanja;
    \item usporedba originalne i poboljšane verzije;
    \item napredna harmonijska i kreativna obrada.
\end{enumerate}

\section{Kriterij uspjeha prve verzije}

Prva verzija je uspješna kada može učitati valjanu MIDI datoteku, prikazati
točan sažetak njezine strukture i glazbenih događaja te završiti bez ikakve
promjene originalne datoteke.

\section{Grafičko sučelje i pokriće aplikacije}

Aplikacija uključuje moderno tamno desktop sučelje s uvozom datoteka
\texttt{.mid} i \texttt{.midi}. Nakon uvoza prikazuje sažetak formata, broj
izvornih trackova, trajanje, početni tempo, ukupan broj nota i status analize.

Glavni dio sučelja sadrži 16 stalnih track-kartica. Svaka kartica prikazuje
broj tracka, izvor podataka, službeni broj i naziv prepoznatog instrumenta,
ikonu procijenjenog instrumenta, broj nota, najveću polifoniju i postotak
pouzdanosti zaključka. Odabirom kartice korisnik dobiva detalje o kanalima,
rasponu nota i dokazima korištenim pri prepoznavanju.

\subsection{Pravila ponašanja aplikacije}

Sučelje je upravljački sloj, a ne mjesto na kojem se skrivaju glazbena pravila.
Sva pravila moraju postojati u testabilnim modulima i davati isti rezultat u
GUI-u, naredbenom retku i automatskom testu.

\begin{enumerate}[label=A\arabic*.]
    \item Uvoz nikada automatski ne pokreće izmjenu. Nakon učitavanja aplikacija
    ulazi u način \emph{Analyze}.

    \item Aplikacija mora razlikovati obični SMF, potvrđeni Pa800 Style MIDI,
    Style Works izvoz, zbirku datoteka i nepodržani format. Odabir parsera mora
    biti prikazan korisniku zajedno s razlogom.

    \item Svaki posao ima jedinstveni ID, izvorni SHA-256, vrijeme pokretanja,
    verziju parsera, verziju registra, verziju skupa pravila i aktivnu
    konfiguraciju.

    \item Rezultat parsera je nepromjenjivi model izvora. Moduli prijedloga i
    izmjena rade nad zasebnom radnom kopijom i ne mogu prepisati originalni
    model.

    \item Kartica mora jasno razlikovati fizički track, logički MIDI kanal,
    Style ulogu i CV. Broj kartice nikada se ne smije predstaviti kao MIDI
    kanal bez oznake.

    \item Uz svaki instrument prikazuju se puni broj, službeni naziv, status
    dokaza i izvor. Kod statusa \texttt{INFERRED} prikazuju se postotak i sva
    pravila koja su sudjelovala u procjeni.

    \item Status \texttt{UNKNOWN} ne smije se automatski tumačiti kao User
    Sound. Aplikacija odvojeno prikazuje \path{FACTORY_CONFIRMED},
    \path{USER_SLOT_CONFIRMED}, \path{UNKNOWN}, \path{CONFLICT} i
    \path{INCOMPLETE_ADDRESS}. User slot potvrđuje lokaciju, ali ne i naziv
    ili stvarni sadržaj zvuka.

    \item Korisnik mora moći odabrati vremenski segment instrumenta i zamijeniti
    ga potvrđenim Factory Soundom. Zamjena se ne primjenjuje samim odabirom u
    pregledniku: GUI zasebno prikazuje ciljnu adresu, pogođeni kanal, početni
    tick, završni tick i gumb \emph{Apply replacement}.

    \item Preglednik Factory Soundova koristi produkcijski registar. Filtrira
    po obitelji instrumenta, glazbenoj ulozi, rasponu i kompatibilnosti
    uređaja. Prijedlog nije automatska odluka; konačni Factory Sound bira
    korisnik.

    \item Zamjena se izvršava po Program Change segmentu, a ne slijepo za cijeli
    fizički track. Ako track mijenja instrument tijekom pjesme, svaki segment
    ima zasebnu karticu, odluku i dnevnik promjene.

    \item Primjena Factory Sounda smije promijeniti samo odobrene Bank Select
    \texttt{CC00}, \texttt{CC32} i Program Change vrijednosti za odabrani
    segment. Note, velocity, tempo, metar, kontroleri, SysEx, tekstovi i drugi
    segmenti ostaju nepromijenjeni. Nedostajući Bank Select ne smije se tiho
    pretpostaviti niti umetnuti bez jasnog prijedloga i potvrde.

    \item Prije primjene prikazuje se razlika \emph{izvorna adresa $\to$ ciljna
    Factory adresa}. Nakon primjene dostupni su Undo, vraćanje izvornog Sounda
    i verifier koji ponovno parsira izlaz i blokira svaku neočekivanu razliku.

    \item Upozorenja imaju razine \texttt{INFO}, \texttt{WARNING},
    \texttt{ERROR} i \texttt{BLOCKING}. \texttt{BLOCKING} onemogućuje Enhance
    i spremanje izmijenjene datoteke, ali ne mora onemogućiti read-only izvoz
    analize.

    \item Korisnik mora moći filtrirati rezultat po Styleu, elementu, CV-u,
    ulozi, kanalu, instrumentu, noti, statusu dokaza i upozorenju bez promjene
    samih podataka.

    \item Prikaz instrumenta mora odvojiti službenu definiciju iz User's
    Manuala od empirijskih kartica \emph{Chordal}, \emph{Line/Riff},
    \emph{Solo candidate}, \emph{Fixed Intro/Ending} i \emph{Guitar/RX}.
    Factory Style statistika ne smije biti prikazana kao tvornička Sound
    postavka niti nuditi automatsku promjenu Sounda.

    \item Na razini svakog Elementa/CV-a aplikacija mora prikazati klasifikaciju
    tracka prije prikaza bilo kakvog velocity prijedloga. Ako klasifikacija nije
    \texttt{ACCOMP\_CHORDAL} ili potvrđena sekundarna perkusija, pravila 35\% i
    50\% ostaju zaključana.

    \item Svaki prijedlog prikazuje staru vrijednost, novu vrijednost, formulu,
    pogođene događaje, mogući rizik i pravilo. Grupna potvrda dopuštena je samo
    za istu vrstu pravila i jasno prikazan opseg.

    \item \emph{Undo} prije spremanja vraća radnu kopiju na prethodno stanje.
    Nakon spremanja povrat se osigurava zadržavanjem originala i dnevnikom
    promjena, a ne prepisivanjem izvora.

    \item Prekid ili greška analize ne smije ostaviti djelomični JSON označen
    kao potpun. Izvještaj dobiva status \texttt{PARTIAL} i popis nedovršenih
    koraka.

    \item Analiza iste datoteke s jednakim pravilima i konfiguracijom mora biti
    deterministička. Ako rezultat ovisi o nasumičnosti, seed mora biti zapisan,
    ali nasumičnost nije dopuštena u osnovnom Analyze načinu.

    \item Aplikacija ne smije slati MIDI sadržaj, nazive ili kontrolne zbrojeve
    na mrežu bez zasebne izričite korisničke dozvole.

    \item Zatvaranje aplikacije tijekom neizvezenog rada mora ponuditi
    odustajanje, spremanje izvještaja ili nastavak rada. Nikada ne smije potajno
    spremiti izmijenjeni MIDI preko originala.
\end{enumerate}

\subsection{Mapiranje MIDI formata}

MIDI Format 0 sadrži jedan fizički track, ali može koristiti svih 16 MIDI
kanala. Aplikacija zato kanale 1--16 prikazuje kao zasebne logičke trackove.
Time se različiti instrumenti unutar jedne Format 0 datoteke ne spajaju u jednu
nejasnu karticu.

MIDI Format 1 sadrži više fizičkih trackova. Aplikacija izravno mapira fizičke
trackove 1--16 na odgovarajuće kartice. Track koji sadrži tempo, oznaku takta i
druge meta-događaje, ali nema note, označava se kao \emph{Conductor / Meta}.
Neiskorišteni slotovi ostaju jasno označeni kao prazni.

\subsection{Prepoznavanje instrumenta i uloge tracka}

Zaključak o instrumentu ne temelji se na samo jednom podatku. Aplikacija
kombinira sljedeće metode:

\begin{enumerate}
    \item \textbf{Program Change} poruku i svih 128 General MIDI programa;
    \item standardni percussion kanal 10 za prepoznavanje bubnjeva i udaraljki;
    \item naziv tracka i ključne riječi poput piano, bass, guitar, strings,
    brass, lead, pad i drums;
    \item prosječni i ukupni tonski raspon;
    \item najveću polifoniju i odnos melodijskog i akordskog sviranja;
    \item profil trajanja nota za razlikovanje melodijske, ritmičke i prateće
    uloge.
\end{enumerate}

Eksplicitna MIDI informacija ima veću težinu od heurističke procjene. Ako
Program Change nije zapisan, aplikacija može predložiti glazbenu ulogu, ali
prikazuje nižu pouzdanost i navodi korištene dokaze. Time se procjena nikada ne
predstavlja kao sigurna činjenica.

\subsection{Registracija svih Pa800 instrumenata i Drum Kitova}

Svaki Pa800 instrument, zvuk, Drum Kit i percussion kit mora bez izuzetka biti
registriran pod svojim punim službenim brojem iz priručnika. Ovo pravilo vrijedi
za sve tvorničke banke i sve kategorije; nijedna stavka ne smije biti
preskočena. Službeni broj sastoji se od MIDI vrijednosti \texttt{CC00.CC32.PC},
odnosno Bank Select MSB, Bank Select LSB i Program Change. Svaka komponenta mora
biti u rasponu 0--127. Broj se prikazuje kao jedna adresa i glavni je
identifikator stavke. Neposredno nakon broja zapisuje se službeni naziv iz
PDF-a, primjerice \texttt{121.0.33 --- Finger Bass GM} ili
\texttt{120.0.5 --- Standard Kit RX1}. Redoslijed je uvijek broj pa naziv.

Broj i naziv moraju se prepisati točno iz službenog PDF-a, bez pretvaranja u
General MIDI broj, skraćivanja, prevođenja ili nagađanja značenja pojedinih
dijelova broja. Instrument se ne smije proglasiti registriranim ako nedostaje
broj ili ako veza između broja i naziva nije potvrđena u priručniku. Ako više
instrumenata ili Drum Kitova ima jednak ili sličan naziv, razlikuju se prema
punom službenom broju. Svaki instrument i svaki Drum Kit mora imati vlastiti
zaseban zapis; jedan zapis ili broj ne smije se zajednički dodijeliti većem broju
različitih stavki.

Registar za svaki instrument i svaki Drum Kit mora sadržavati najmanje puni
službeni broj, službeni naziv, vrstu stavke i broj stranice PDF-a na kojoj je
par potvrđen. Sučelje, tekstualni sažetak i JSON izvještaj prikazuju službeni
broj prije naziva. Potpunost registra provjerava se usporedbom broja zapisa sa
svim stavkama navedenim u službenom PDF-u. Ako nedostaje makar jedan instrument
ili Drum Kit, registar ne može dobiti status \texttt{PASS}. Ako službeni PDF
nije dostupan ili ne prođe provjeru kontrolnog zbroja, Pa800 registar dobiva
status \texttt{BLOCKED} i ne popunjava se prema neslužbenim izvorima.

Službene bankovne tablice na tiskanim stranicama 275--283 i zasebna Drum Kit
tablica na tiskanoj stranici 295 izvor su produkcijskog K01 registra.
Reproducibilni generator prihvaća samo potvrđeni SHA-256 priručnika i pinned
PDF extractor, stvara kanonski JSON, a produkcijski lookup radi isključivo za
čitanje. Aktualni generator i testovi potvrđuju 1.006 Factory Sound adresa i 65
izričito imenovanih Drum Kit adresa, ukupno 1.071 jedinstvenu punu adresu.
Adrese 120.0.48 i 120.0.49 sačuvane su kao dvije zasebne adrese sa službenim
nazivom \texttt{Orchestra Kit GM}. K01 ima ograničeni read-only
\texttt{PASS}; ne obuhvaća K02 rezoluciju niti MIDI izmjene.

Na tiskanoj stranici 295 priručnik istodobno navodi imenovane stavke
\texttt{120.0.57 --- SFX Kit 2} i \texttt{120.0.58 --- Synth Kit} te raspon
\texttt{57--63 (remap to 56)}. K01 čuva izričito imenovane stavke 57 i 58,
označava preklop kao \texttt{CONFLICT} i ne primjenjuje remap preko imenovanih
stavki. Nepreklopljeni dio 59--63 ostaje zabilježen u službenom remap retku za
budući K02, bez promjene MIDI događaja.

\subsection{Što aplikacija može detektirati}

Za cijelu MIDI datoteku aplikacija detektira format, rezoluciju, trajanje,
tempo i njegove promjene, oznake takta i tonaliteta, broj trackova, kanale,
programe, kontrolne poruke, sustain pedal, pitch bend, percussion događaje i
strukturne probleme.

Za svaki track ili logički kanal detektira naziv, broj nota, tonski raspon,
velocity statistiku, gustoću nota, prosječno trajanje, najveću polifoniju,
omjer monofonih i akordskih početaka te korištene General MIDI programe. Na
temelju tih mjerenja može prepoznati ili predložiti sljedeće uloge:

\begin{itemize}
    \item conductor ili meta track;
    \item drum i percussion groove;
    \item bass liniju;
    \item lead ili melodijsku liniju;
    \item akordsku pratnju;
    \item pad ili dugotrajnu harmonijsku teksturu;
    \item arpeggio ili brzi motiv;
    \item rhythm guitar, solo guitar i power-chord guitar;
    \item opći melodijski ili ritmički part kada nema dovoljno dokaza za užu
    klasifikaciju.
\end{itemize}

\subsection{Pravila za prepoznavanje glazbene uloge}

Pravila su deterministička i prikazuju korištena mjerenja. Eksplicitni podaci,
poput MIDI kanala i Program Change poruke, uvijek imaju prednost pred
heuristikom.

\begin{enumerate}[label=D\arabic*.]
    \item \textbf{Percussion} se potvrđuje kada se najmanje 50\% nota nalazi na
    standardnom MIDI kanalu 10.

    \item \textbf{Bass line} se predlaže kada je prosječna visina niža od MIDI
    note 48 i najmanje 65\% početaka nota je monofono.

    \item \textbf{Pad / sustained harmony} zahtijeva prosječno trajanje od
    najmanje dvije četvrtinke i najveću polifoniju od najmanje tri note.

    \item \textbf{Chordal accompaniment} se predlaže kada najmanje 35\%
    početaka sadrži dvije ili više istodobnih nota ili kada je najveća
    polifonija najmanje četiri.

    \item \textbf{Arpeggio / fast motif} zahtijeva najmanje dvije note po
    četvrtinki, najmanje 75\% monofonih početaka i prosječno trajanje note do
    0,75 četvrtinke.

    \item \textbf{Lead / melody} zahtijeva prosječnu visinu od najmanje MIDI
    note 60 i najmanje 75\% monofonih početaka.

    \item Ako nijedno pravilo nema dovoljno dokaza, koristi se neutralna oznaka
    \emph{Melodic / rhythmic part} s nižom pouzdanošću.
\end{enumerate}

\subsection{Posebna pravila za gitaru}

Guitar uloga analizira se samo kada Program Change pripada General MIDI guitar
programima 25--32 ili kada naziv tracka izričito sadrži riječ \emph{guitar} ili
\emph{gitara}. Time se sprječava proglašavanje svakog akordskog instrumenta
gitarom.

\begin{enumerate}[label=G\arabic*.]
    \item \textbf{Power-chord guitar} zahtijeva najmanje dva akordska udara u
    kojima postoji odnos root--fifth, odnosno interval čiste kvinte. Struktura
    ne smije sadržavati malu ili veliku tercu, a najmanje 50\% svih akordskih
    udara mora zadovoljiti ovo pravilo. Oktava je dopuštena.

    \item \textbf{Solo guitar} zahtijeva najmanje 75\% monofonih početaka,
    najveću polifoniju do dvije note i najviše 25\% akordskih početaka. Pitch
    bend nije obavezan, ali povećava pouzdanost solo zaključka.

    \item \textbf{Rhythm guitar} prepoznaje se kada najmanje 30\% početaka
    sadrži akord ili kada najveća polifonija iznosi najmanje tri. Pravilo za
    power chord provjerava se prvo, kako power-chord part ne bi bio označen samo
    kao opća rhythm guitar.

    \item Ako je guitar instrument potvrđen, ali nijedan obrazac nije dovoljno
    izražen, prikazuje se \emph{Guitar part (uncertain role)} umjesto izmišljene
    precizne klasifikacije.
\end{enumerate}

Ova pravila opisuju sadržaj MIDI događaja, a ne stvarni zvuk snimke. Ako MIDI
ne sadrži ispravan Program Change ili naziv tracka, aplikacija može pouzdano
opisati obrazac nota, ali ne može sa sigurnošću dokazati da je part odsviran na
gitari.

\subsection{Ikone instrumenata}

Kartice koriste ugrađene glazbene ikone za piano, gitaru, bass, gudače, brass,
reed, flaute, synthesizer, vokal, efekte i percussion. Uz ikonu se uvijek
prikazuje tekstualni naziv, pa sučelje ostaje razumljivo i na sustavu koji nema
potpunu podršku za slikovne glifove.

\subsection{Rezultati i sigurnost}

GUI radi u načinu \emph{Analyze}: ne mijenja uvezeni MIDI. Analiza provjerava
SHA-256 kontrolni zbroj originala, prikazuje upozorenja te omogućava izvoz
potpunog JSON izvještaja. Dodatne biblioteke koriste se automatski kada su
dostupne, dok ugrađeni parser osigurava osnovni rad bez vanjskih MIDI paketa.

\section{Status prve verzije}

U trenutačnom snapshotu implementirani su učitavanje i validacija MIDI
datoteke, analiza glazbenih događaja, provjera očuvanja originala, tekstualni i
JSON izvještaj te automatsko povezivanje dostupnih biblioteka. Ugrađeni parser
ostaje osnovni mehanizam, \texttt{numpy} proširuje statistiku, a
\texttt{mido}, \texttt{pretty\_midi} i \texttt{music21} uključuju se kada su
dostupni. Mapiranje formata 0 i 1, prikazni model od 16 trackova, višemetodno
prepoznavanje instrumenata te klasifikacija rhythm, solo i power-chord guitar
uloga imaju opće MIDI testove. Format 2 sada ima zaseban read-only testni scope:
svaki track ostaje neovisna sekvenca, bez izmišljene globalne tempo mape ili
trajanja. \texttt{tkinter} 8.6 je dostupan, a workspace-only Xvfb omogućio je
stvarni automatizirani test Tk prozora, 16 kartica, Format 0 analize/JSON izvoza
i Format 2 read-only prikaza. Native desktop pakiranje i ljudski UX ostaju
necertificirani. Produkcijski K01 registar, generator, kanonski adresni JSON,
immutable lookup i namjenski testovi sada su prisutni; K01 ima ograničeni
read-only \texttt{PASS}. K02 resolver sada povezuje stvarne M10 segmente s K01
dokazom i ima ograničeni read-only \texttt{PASS}. Immutable Change Plan i
Suggest Policy Engine također imaju ograničeni \texttt{PASS}, ali uvijek vraćaju
\texttt{apply\_authorized=false}. K03-P user-selected proposal builder,
in-memory Change Engine, rollback, independent Verifier i atomic verified
writer sada su prisutni u ograničenom scopeu. Bounded programatska K03 zamjena
smije stvoriti samo novi \texttt{\_enhanced} MIDI i JSON nakon svih USER/V01 kapija.
Bounded GUI workflow za izbor, preview, USER potvrdu, V01, W01 save i rollback
također je prisutan; automatski Enhance nije uključen. Ostali
Pa800-specifični moduli dobivaju zaseban status tek nakon vlastite
implementacije i dokaza. Pravila za buduće načine \emph{Suggest} i
\emph{Enhance} ostaju obavezna pri nastavku razvoja.

Read-only Style Works Loader ima status \texttt{PASS}. Automatski je učitao i
validirao svih 3.211 MIDI datoteka iz arhiva, prepoznao 252 Style skupa,
razdvojio Element/CV/ulogu i dokazao očuvanje izvornog ZIP SHA-256 zbroja.
Element \path{Fox Shuffle 1_Break.mid} s konfliktnim oznakama \texttt{1 Bar} i
\texttt{5 Bars} sačuvan je, ali ispravno blokiran za vremensku analizu.

Read-only Context Classifier M07 ima status \texttt{PASS}. Generička jezgra i
Style adapter odvojeno vraćaju strukturnu ulogu, glazbenu funkciju, encoding,
fixed Intro/Ending kandidat, RX rizik i sigurnosnu politiku. Drum, Perc i Bass
imaju prioritet pred melodijskom heuristikom; bliski onseti grupiraju se prema
PPQ-u; Guitar Mode se ne zaključuje samo iz raspona nota; višekanalni konflikt
blokira klasifikaciju. Puni corpus smoke test obradio je 3.211 datoteka i
65.021 track-slice rezultat bez iznimke, uz nepromijenjen SHA-256 arhiva.

M07 je isključivo analizator i ne autorizira \emph{Enhance}.

Analyze-only Measure/Phrase Analyzer M08 ima status \texttt{PASS}. Iz
eksplicitnih tempo i metar događaja gradi vremensku mapu, određuje mjeru, beat
i tick svakog događaja, pronalazi egzaktno ponovljene mjere te označava
konzervativne kandidate granica fraza nakon duge stanke. Konfliktne oznake
tempa ili metra blokiraju vremensko tumačenje, a Style adapter koristi samo
potvrđeni M06 vremenski prozor. M08 ne kvantizira, ne pomiče i ne briše MIDI
događaje te ne autorizira \emph{Humanize}, strumming ili drugi \emph{Enhance}.

Read-only Instrument Measurement Engine M09 ima status \texttt{PASS}. Za svaki
potvrđeni part računa raspon nota, velocity, trajanje u tickovima i beatovima,
gustoću, sounding i onset polifoniju, kontrolere, pitch bend i aftertouch. Kada
su dostupni M08 podaci, stvara zasebne statistike po mjeri i kandidatu fraze.
Nezatvorene ili orphan note daju status \texttt{PARTIAL}, dok blokirani M07 ili
M08 kontekst daje \texttt{BLOCKED}. M09 nasljeđuje najstrožu sigurnosnu
politiku i ne sadrži \emph{Suggest}, \emph{Change}, \emph{Humanize} ili writer
funkciju.

Read-only Instrument Segmenter M10 ima status \texttt{PASS}. Po MIDI kanalu
prati \texttt{CC00}, \texttt{CC32} i Program Change te razdvaja part na
vremenske segmente sa zasebnim M09 profilom. Potpuna adresa prikazuje se kao
\texttt{CC00.CC32.PC}; nepotpuna ostaje \texttt{INCOMPLETE}, a događaji prije
prvog programa \texttt{NO\_PROGRAM}. Izvorni event order određuje granicu kada
događaji dijele tick, a note preko Program Change granice dobivaju upozorenje.
M10 ne proglašava Factory/User identitet: svaki segment ostaje
\texttt{UNRESOLVED} i nijedan se događaj ne mijenja. Datoteka
\texttt{instrument\_segmenter.py} i njezini namjenski testovi prisutni su u
trenutačnom Git snapshotu. K01 Factory lookup i K02 read-only resolver sada su
dovršeni. K02 vraća zaseban immutable rezultat i namjerno ne prepisuje izvorni
M10 \texttt{UNRESOLVED} segment niti bilo koji MIDI događaj.

Naknadni neovisni testni audit proširio je izravno pokriće M08--M10. Svaki od
ta tri modula sada ima devet namjenskih testova. Ispravljeni su potpuni potpis
egzaktnog ponavljanja, faza mjere u prozoru koji ne počinje na ticku nula,
provjera SHA-256/PPQ/prozora M09 timelinea, višetrack zaštita, gustoća bez
potvrđenog prozora, otvorena nota preko Program Change granice, jedinstveni
ordinali i neodređeni cross-track redoslijed.

Zaseban izvršni corpus test ponovno obrađuje svih 3.211 Factory MIDI datoteka
i 65.021 track-slice. Provjerava 82.917 M10 segmenata, fiksne statusne
raspodjele, deterministički rezultatni digest i nepromijenjen SHA-256 izvornog
ZIP-a. Status \texttt{PASS} ostaje ograničen na izravno testirani read-only
scope; GUI, identity, writer, Enhance i hardware ponašanje time nisu
certificirani.

Evidence Registry Foundation G00 ima status \texttt{PASS}. Kanonski JSON
registar i deterministički generirani Markdown odvajaju dokaz, implementaciju,
validaciju i certifikaciju. Validator provjerava ID namespace, reference,
lokalne putanje, SHA-256 zbrojeve, prisutnost registriranih testova i obaveznu
izolaciju dokumenata označenih kao \texttt{IDEA\_CATALOG}. Time se ne tvrdi da
validator sam dokazuje glazbenu ispravnost: svaka domenska tvrdnja i dalje
zahtijeva vlastiti reproducibilni test i izvor.

M07 corpus rezultat sada ima zaseban reproducibilni test koji ponovno obrađuje
svih 3.211 datoteka i provjerava 65.021 rezultat, uključujući raspodjelu
\texttt{CLASSIFIED}, \texttt{UNCERTAIN} i \texttt{BLOCKED}. DNA paket je
zabilježen kao djelomično validiran jer sadrži i bajtovno jednaku kopiju
Factory arhiva; ta kopija ne smije ući u budući Gold training ili evaluation
skup.

\section{Završno potvrđeno stanje i sljedeći razvojni korak}

U stvarnom snapshotu prisutni su opći MIDI moduli M01--M05, read-only moduli
M06--M10 i G00 Evidence Registry Foundation. Aktualni puni testni run mora se
zapisati u \texttt{MODULE\_LOG.md}; povijesne brojke same nisu dokaz. Navedeni
read-only moduli ne autoriziraju \emph{Enhance}.

K01 Factory registry ima produkcijski generator, kanonski JSON, immutable
lookup i namjenski test te ograničeni read-only status \texttt{PASS}. K02
resolver nad stvarnim M10 segmentima također ima ograničeni read-only
\texttt{PASS}. Immutable Change Plan, Suggest Policy Engine, K03-P proposal,
in-memory Change Engine, independent Verifier i atomic writer sada su prisutni.
Bounded K03 backend i automatizirani GUI save/rollback workflow imaju
ograničeni \texttt{PASS}; automatski \emph{Enhance} ostaje \texttt{NO-GO}.
Implementirani K02 statusi su
\texttt{FACTORY\_CONFIRMED}, \texttt{USER\_SLOT\_CONFIRMED},
\texttt{UNKNOWN}, \texttt{CONFLICT}, \texttt{INCOMPLETE\_ADDRESS} i
\texttt{NO\_PROGRAM}; nijedan status sam ne autorizira promjenu.

A01 Workspace Persistence Probe i A02 Repository Inventory Auditor imaju
ograničeni governance \texttt{PASS}. A01 checkpoint ponovno je provjeren nakon
sljedeće korisničke poruke: commit, Python datoteke i testovi ostali su prisutni
i bajtovno jednaki. A02 inventar nad aktualnim \texttt{git ls-tree} popisom
nema error ni warning nalaz za dokumentirane produkcijske, testne i generatorske
artefakte. Ovaj status ne certificira glazbenu ispravnost. M01 Format 2 test
sada je dovršen za ograničeni read-only scope. Izravni core-parity guard također
uspoređuje oba parsera na sintetičkim rubovima i svih 3.211 Factory MIDI
ulaza. Automatizirani Tk runtime također je potvrđen na lokalnom virtualnom
displayu. Reprodukcijski K01 generator, lookup, K02 resolver, immutable Change
Plan, Suggest Policy Engine i K03-P user-selected proposal builder sada su
dovršeni u svojim ograničenim scopeovima. Windows \texttt{install.bat} i
\texttt{run.bat} imaju statički testirani status \texttt{PARTIAL}; stvarno
Windows/Tk izvršavanje čeka korisnički računar. Bounded GUI tok za USER
selection, preview, approval, independent verification, atomic save i rollback
sada je dovršen u automatiziranom virtual-display scopeu. R01 Factory
kontekstualni profili također su dovršeni kao reference-only podaci bez
miješanja uloga. S01 exact-profile velocity outlier pravilo sada može stvoriti
samo immutable Suggest plan uz stroge context/support kapije; C01 namjerno ne
podržava primjenu Note On velocity mutacije. G02 sada isključuje exact Factory
container iz DNA paketa, deduplicira Gold i potvrđuje source-disjoint evaluation
manifest. M07 statistička kalibracija ostaje \texttt{BLOCKED} jer Gold nema
ručno potvrđene per-part role labele. V02 specialized velocity Apply/Verifier
scope sada postoji, ali zahtijeva zasebni USER acknowledgement da su
velocity-switch/RX pragovi nepoznati te provodi max 8/max 32/1--127 kapije.
Generic i automatski Apply ostaju blokirani. Sljedeći korak je povezati S01/V02
s GUI preview/approval/risk-ack/save tokom; automatski Enhance ostaje zabranjen.

\end{document}
